\note{sa page « 5 bis » ; le paragraphe 5 est écrit après coup et inséré ici}

5. On peut aussi donner une condition d'existence des $\pi_{i}$, liée aux
conditions (4.4). Pour ceci, on suppose que $\mathcal{E}$ est contenu dans un
sur-anneau gradué $\mathcal{F}$ \add{p.\ ex.\ l'anneau des
endom.\ d'un $M^{*}$ comme dans 1.} pour lequel les $\pi_{i}$ satisfont (1.1.)
et les relations correspondantes (1.3) (avec $\mathcal{E}$ remplacé par
$\mathcal{F}$) \add{et $(3.0)$ : $\pi_{i} = 0$ si $i \notin [0,2n]$} sont
satisfaites. \struck{Supposons} Supposons que pour tout $i$ la condition
$(4.4)_{i}$ soit satisfaite (avec $v_{i}$, $w_{i} \in \mathcal{E}$ bien sûr).
On en conclura que

\[ (5.1) \qquad \pi_{i} \in \mathcal{E} \quad \text{pour tout } i . \]

\struck{Raisonnons} \struck{Supposons provisoirement} \ill{} $\pi_{\alpha} \in
\mathcal{E}$ si $\alpha \notin [i, 2n-i]$, \uncertain{l'on veut} vérifier \ill{}
\struck{si $i \leqslant n$,} i.e.\ $\pi_{i}$ et $\pi_{2n-i} \in \mathcal{E}$.
(Trivial si $i = n-1$.) \struck{Supposons pour $i \leqslant n-2$}
\add{comme $\varphi_{i}$, $\psi_{2n-i} \in \mathcal{E}$, et} pour $i+1$,

Notons que, avec les notations de 3., \struck{\ill{}}
$\psi_{2n-i} v_{i} \varphi_{i}$, $\varphi_{i} w_{i} \psi_{2n-i} \in
\mathcal{E}$, i.e.
\[ v'_{i} = \psi_{2n-i}\, v_{i}\, \varphi_{i} = \pi_{2n-i}\, v_{i}\, \pi_{i} \in \mathcal{E} , \]
\[ w'_{i} = \varphi_{i}\, w_{i}\, \psi_{2n-i} = \pi_{i}\, w_{i}\, \pi_{2n-i} \in \mathcal{E} , \]
et \struck{comme}
\[ w'_{i} v'_{i} = \pi_{i}\, w_{i}\, \pi_{2n-i}\, v_{i}\, \pi_{i} = \pi_{i}\, \pi_{i}\, w_{i}\, v_{i}\, \pi_{i} = \pi_{i}^{3} = \pi_{i} \]
\marginal{car $w_{i} \pi_{2n-i} = \pi_{i} w_{i}$ ; car $w_{i} v_{i} \pi_{i} = \pi_{i}$}
et de même $v'_{i} w'_{i} = \pi_{2n-i}$, on conclut que $\pi_{i} \in
\mathcal{E}$, $\pi_{2n-i} \in \mathcal{E}$.

